% SPDX-License-Identifier: CC-BY-SA-4.0
% Author: Matthieu Perrin

\begingroup

\SetKwFunction{Read}{read}
\SetKwFunction{Write}{write}
\SetKwData{Value}{value}
\SetKwData{Time}{time}
\SetKwData{Writer}{writer}

\SetKwBroadcast{MB}{mb}
\SetKwMessage{Synch}{synch}
\SetKwMessage{WRITE}{write}
\SetKwVariable{TS}{ts}
\SetKwVariable{Val}{val}

\begin{exercice}[Implémentation d'un registre séquentiellement cohérent]

  On considère l'algorithme d'Attiya, Bar-Noy et Dolev (algo.~\ref{algo:ABD}),
  qui implémente un registre linéarisable.
  On cherche dans cet exercice à l'adapter pour garantir seulement une propriété plus faible :
  la cohérence séquentielle (déf.~\ref{def:SC}).

  \begin{algorithm}[H]
    \lLVariables{\At $p_i$}{
      $\Value_i \leftarrow \bot$;
      $\Time_i \leftarrow \langle 0, 0 \rangle$
    }
    \Method{$\Read()$ \At $p_i$}{
      \nl \MB.\SBroadcast $\Synch()$\;
      \nl \Let $\Val \leftarrow \Value_i$\;
      \nl \MB.\SBroadcast $\WRITE(\Val, \Time_i)$\;
      \nl \Return $\Val$
    }
    \Method{$\Write(\Val)$ \At $p_i$}{
      \nl \MB.\SBroadcast $\Synch()$\;
      \nl \Let $t_i$ \St $\langle t_i,i \rangle > \Time_i$\;
      \nl \MB.\SBroadcast $\WRITE(\Val, \langle t_i,i \rangle)$;
    }
    \Upon{\MB.\Deliver $\WRITE(\Val, \TS)$ \From $p_i$ \At $p_j$}{
      \nl \lIf{$\Time_j <_{lex} \TS$}{
        $\Value_j \leftarrow \Val$;~ $\Time_j \leftarrow \TS$
      }
    }
    \caption{Implémentation d'un registre linéarisable}
    \label{algo:ABD}
  \end{algorithm}

  \begin{definition}[Cohérence séquentielle]\label{def:SC}
    Une exécution $H$ est \emph{séquentiellement cohérente} s'il existe un ordre total
    sur ses opérations qui respecte l'ordre séquentiel des opérations de chaque processus,
    et tel que l'exécution des opérations dans cet ordre retourne les mêmes valeurs que dans $H$.
  \end{definition}

  \begin{questions}
  \item Simplifier autant que possible l'algorithme~\ref{algo:ABD}, tout en garantissant
    que le registre obtenu reste séquentiellement cohérent.
    Justifier que l'algorithme obtenu vérifie la définition~\ref{def:SC}.
    
  \item Considérez deux registres distincts, chacun implémenté par une instance de l'algorithme
    obtenu à la question précédente.
    Construisez une exécution dans laquelle chacune des deux instances est séquentiellement cohérente,
    mais où l'exécution de l'ensemble des opérations sur les deux registres ne l'est pas.
  \end{questions}

\end{exercice}

\endgroup
\endinput
